\section{Variables recorded by only one channel: the visible failure}\label{sec:channel}

A covariate recorded in only one channel breaks support even if the units are comparable. If it does not predict the outcome given the others, dropping it restores identification; if it confounds, neither using it nor dropping it identifies $\theta$.

\paragraph{Measurement model.} Each unit has latent covariates $X^*=(V,C^*)$, where $C^*$ takes values in $\mathcal C\cup\{\vacio\}$ and $C^*=\vacio$ means ``no actual data''. What is recorded depends on the channel: $X=(V,C)$, with $C=C^*$ if $D=1$ and $C=\vacio$ if $D=0$. Section~\ref{sec:instruments} treats the general case $X=\varphi_D(X^*)$. Examples: a commercial history kept only by premium-channel advisors, or diagnostic tests performed only in hospital. We write $P^*$ and $Q^*$ for the laws of $X^*$ in the target group and in the donor channel.

\begin{lemma}[Singular mass]\label{lem:singular}
$\varepsilon=\Ps(\X)=\sup\{P(B):\ B\in\A,\ Q(B)=0\}$.
\end{lemma}
\begin{proof}
Let $N$ be a $Q$-null set on which $\Ps$ is concentrated. For every $B$ with $Q(B)=0$, $\Pac(B)=0$ because $\Pac\ll Q$; hence $P(B)=\Ps(B)\le\Ps(\X)$. Moreover $P(N)=\Pac(N)+\Ps(N)=\Ps(\X)$, so the supremum is attained at $N$.
\end{proof}

\begin{proposition}[A channel variable breaks support]\label{prop:breaks}
Under the measurement model, let $\kappa=\PP(C^*\ne\vacio\mid S=1)$ be the fraction of the target group with actual data. Then $\varepsilon\ge\kappa$, and the ideal classifier equals 1 $M$-almost surely on the profiles with $C\ne\vacio$. If $\kappa>0$, (A3) fails regardless of how the units compare.
\end{proposition}
\begin{proof}
Let $B=\mathcal V\times\mathcal C$, the set of profiles with recorded data. By the measurement model, $Q(B)=\PP(C\ne\vacio\mid D=0)=0$, whereas $P(B)=\kappa$; Lemma~\ref{lem:singular} gives $\varepsilon\ge\kappa$. Since $\PP(Z=0,\tilde X\in B)=(1-\pi)Q(B)=0$, $p=1$ $M$-almost surely on $B$.
\end{proof}

Throughout, $g$ denotes a version of $\E[Y(0)\mid V,D=0]$, and $P_V$ and $Q_V$ the laws of $V$ in the target group and in the donor channel.

\begin{proposition}[Dropping a variable that does not predict the outcome]\label{prop:drop}
Assume (A1), (A4) and:
\begin{enumerate}
  \item \emph{Latent support:} $P^*\ll Q^*$.
  \item \emph{Latent (A2):} $\E[Y(0)\mid X^*,S=1]=\E[Y(0)\mid X^*,D=0]$, $P^*$-almost surely.
  \item \emph{$C^*$ does not predict the outcome in the basic channel given $V$:} $\E[Y(0)\mid X^*,D=0]=\E[Y(0)\mid V,D=0]$, $Q^*$-almost surely.
\end{enumerate}
Then, with $V$ as the only covariates, (a) $P_V\ll Q_V$ and (b) $\E[Y(0)\mid V,S=1]=\E[Y(0)\mid V,D=0]$ $P_V$-almost surely, so Theorems~\ref{thm:identification}--\ref{thm:lebesgue} apply with $X=V$.
\end{proposition}
\begin{proof}
(a) If $Q_V(A)=0$, then $Q^*(A\times(\mathcal C\cup\{\vacio\}))=0$, and by (1) also $P^*(A\times(\mathcal C\cup\{\vacio\}))=P_V(A)=0$.
(b) By (3), $\E[Y(0)\mid X^*,D=0]=g(V)$ except on a $Q^*$-null set, which by (1) is also $P^*$-null. By (2), $\E[Y(0)\mid X^*,S=1]=g(V)$ $P^*$-almost surely, and since $V$ is a function of $X^*$, the tower property within $\{S=1\}$ gives $\E[Y(0)\mid V,S=1]=g(V)$.
\end{proof}

Condition (3) is an outcome-based covariate selection criterion \citep{deluna2011}: it suffices even if $C^*$ is associated with the channel.

\paragraph{When the variable does confound.} If (1) and (2) hold but (3) fails, let $\mu^*(x^*)=\E[Y(0)\mid X^*=x^*,D=0]$; then $\theta=\int\mu^*\dd P^*$. Using only $V$, the weighting functional equals $\int g\dd P_V=\int g(V)\dd P^*$, with bias
\[
\int g\dd P_V-\theta=-\int_{\X^*}(\mu^*-g)\dd P^*.
\]
Under $Q^*$ this integral is zero by definition of $g$. Under $P^*$ it can be nonzero only if $C^*$ predicts $Y(0)$ given $V$ and its distribution given $V$ differs between groups: the two conditions of the classical definition of a confounder, necessary but not sufficient, because the effects could offset. The exact condition for (A2) with $V$ as covariates is $\E_{P^*}[\mu^*-g\mid V]=0$ $P_V$-almost surely; the bias itself vanishes if and only if $\int(\mu^*-g)\dd P^*=0$, which can also happen by cancellation. (3) is sufficient for both and does not involve $P^*$. When the bias is nonzero, neither option identifies $\theta$: without $C$, (A2) fails, and with $C$, (A3) fails and, as shown next, (A2) on the supported part generally fails too. Hence the bounds of Corollary~\ref{cor:bounds}, which assume (A2) there, do not apply either. What remains is a sensitivity analysis on (3).

\paragraph{Keeping the variable does not fix the supported part.} With $C$ among the covariates, the supported part is $\{C=\vacio\}$. But ``$\vacio$'' does not mean the same in both groups: in the donor channel it mixes units with $C^*=\vacio$ and $C^*\ne\vacio$, because $C$ is never recorded there; in the target group it means $C^*=\vacio$. Since the whole donor channel has $C=\vacio$, $m_0(v,\vacio)=g(v)$, and the weighting functional equals $\E_{P^*}[g(V)\mid C^*=\vacio]$, which in general is neither $\theta$ nor $\thac$. Indeed, (A2) on $\{C=\vacio\}$ requires $\mu^*(V,\vacio)=g(V)$ $P^*$-almost surely on $\{C^*=\vacio\}$, that is, that among donors with the same $V$ those with and without a history respond alike in the basic channel; (3) implies it, but it is weaker. The doubly robust estimator, in addition, extrapolates the outcome model to units with data (Corollary~\ref{cor:drnosupport}). That model is fitted in the donor channel, where $C$ never varies; if it does not depend on $C$ outside the support, as with trees, which cannot split on a variable that is constant in training, it extrapolates $g(V)$, and the doubly robust estimator targets $\E_{P^*}[g(V)]$, as if $C$ were not used.

\paragraph{Testability.} Condition (3) concerns $Y(0)$ in the basic channel, where $C^*$ is not observed, so it cannot be tested directly. That $C$ predicts the outcome in the premium channel is a hint but not proof, because $Y(1)$ is observed there.

\paragraph{Relation to the missingness-pattern approach.} \citet{dagostino2000} and \citet{blake2020} estimate the propensity within each missingness pattern and require the partially observed variable not to confound in the pattern where it is missing. A channel variable is the case in which the treatment determines the pattern: in the observed pattern there are no donor units (Proposition~\ref{prop:breaks}), and in the missing one it requires $\mu^*(V,\vacio)=g(V)$, which (3) implies.

\paragraph{Practical diagnostic.} For the ideal classifier the failure is visible: including a channel variable creates target profiles with $p=1$. For estimated classifiers, removing a channel variable should reduce the AUC and the estimated mass $\hat P(\hat p>1-t)$, and a large jump in that mass, or in the estimated weights on the $\pi$-invariant scale of Section~\ref{sec:classifier}, when a variable is included is the signature of Proposition~\ref{prop:breaks}; Section~\ref{sec:lalonde} shows that this signal can be faint, and that univariate checks by channel can be more informative. In the simulation, the estimated unsupported mass with the channel variable is $0.326$ against $\kappa=0.335$, and the doubly robust estimator behaves as Corollary~\ref{cor:drnosupport} predicts rather than as the weighting estimator does (Section~\ref{sec:res-main}).

\subsection{Illustration: an experimental benchmark with a top-coded covariate}\label{sec:lalonde}

A covariate need not be missing to break support: it suffices that one channel records it with a ceiling. Public data with an experimental benchmark show such a failure, and show that estimated classifiers barely register it.

\paragraph{Data.} We use the subsample of \citet{dehejia1999} of the National Supported Work experiment, with 185 participants and 260 randomized controls, and the CPS-1 comparison group of \citet{lalonde1986}, with 15{,}992 people, in the copies distributed by \citet{cunningham2021}. The target group is the participants and the donors are the CPS sample. The estimand is the mean of their 1978 earnings without the program, whose experimental benchmark is the mean of the controls: 4{,}555 dollars (standard error 340). The covariates are age, education, race, marital status, lack of a degree, earnings in 1974 and 1975 with indicators of zero earnings, and the squares of age and education. This analysis was not registered in advance; we report the median over five cross-fitting partitions.

\paragraph{Results.} The unadjusted CPS mean is 14{,}847 dollars. Weighting alone gives 5{,}092 with a logistic classifier and 3{,}990 with LightGBM \citep{ke2017}, whereas the doubly robust estimator gives 4{,}783 and 4{,}716 (standard errors 397 and 876), both within one standard error of the benchmark (340). Overlap is poor: with LightGBM, $\neff$ is 33 out of 15{,}992 donors.

\paragraph{Top-coding in one source.} The two sources record earnings differently \citep{smith2005}. One difference is visible: CPS earnings are top-coded (19\,\% of 1974 earnings sit exactly at the cap of 25{,}862 dollars) and NSW earnings are not. Two participants (1.1\,\%) earned more than the cap in 1974, so their profiles have no support by construction: a singular mass in the sense of Lemma~\ref{lem:singular}, on which the ideal classifier equals 1. No estimated classifier comes close. The participant with 35{,}040 dollars gets $\hat p$ between 0.008 and 0.012 with the logistic classifier, about the base rate $\pi=0.011$, and at most 0.001 with LightGBM, in all five partitions: the donors at the cap make high earnings look donor-like. The participant only 67 dollars above the cap gets $\hat p$ between 0.17 and 0.33 with the logistic classifier, 18 to 42 times the prior odds, a noticeable but far from conclusive signal, and between 0.002 and 0.20 with LightGBM. Here the failure does not matter: applying the cap to both groups changes the estimates by less than 30 dollars. The outcome is also top-coded: 26\,\% of CPS 1978 earnings sit at the cap, which biases the donor outcome model downward; applying the cap to the controls lowers the benchmark by 54 dollars, to 4{,}501, small relative to the standard errors. But the example shows that classifier-based diagnostics can miss a support failure that a histogram of each covariate by source, with a fifth of the donor values piled at the maximum, reveals at once. Top-coding is in fact a channel-specific instrument, $x\mapsto\min(x,\bar x)$, that breaks support only at its edge. Section~\ref{sec:instruments} turns to instruments that do not break support at all; other differences in how the two sources collected earnings would be of that kind.
